\section{Заключение}

В данной работе был проведен сравнительный анализ аналитических и численных методов решения задачи о движении тела в гравитационном поле. Основные результаты работы:

\begin{enumerate}
    \item Получено аналитическое решение для случая однородного гравитационного поля без учета сопротивления воздуха (формулы \eqref{eq:newton_law}--\eqref{eq:trajectory}).
    \item Разработаны и реализованы на языке \texttt{Python} численные алгоритмы: метод Эйлера и метод Рунге-Кутты 4-го порядка (листинг~\ref{code:python}).
    \item Проведен сравнительный анализ точности различных методов на тестовых примерах (таблица~\ref{tab:comparison}). Показано, что метод RK4 обеспечивает погрешность менее 0.2\% при шаге интегрирования \(h=0.05\) с, в то время как погрешность метода Эйлера при том же шаге достигает 3.5\%.
    \item Исследовано влияние сопротивления воздуха на траекторию движения тела (рисунки~\ref{fig:trajectories} и~\ref{fig:drag}). Показано, что при типичных параметрах дальность полета уменьшается на 15--25\%.
\end{enumerate}

Полученные результаты могут быть использованы при разработке программного обеспечения для баллистических расчетов, в учебном процессе при изучении численных методов, а также в качестве основы для более сложных моделей, учитывающих дополнительные факторы (неоднородность поля, ветер, вращение Земли).